\documentclass[paper=a4, fontsize=11pt]{scrartcl}

\newcommand{\assignment}{Project Proposal Template}
\newcommand{\duedate}{March 30, 2026}
\input{include/hw-template.tex}
\author{
    \textbf{Gin Fu, Shiran Wang} \\
    \href{mailto:fu.jingyi2@northeastern.edu}{fu.jingyi2@northeastern.edu}, \href{mailto:wang.shira@northeastern.edu}{wang.shira@northeastern.edu} \\
    \href{https://github.com/blueif16/dreamrag}{github.com/blueif16/dreamrag}
}

\begin{document}
\maketitle

\section{Proposed Objective}

We propose DreamRAG, a retrieval-augmented dream-analysis system whose
response to each query is a bento-grid dashboard rather than a text reply. A
self-correcting LangGraph agent retrieves over three corpora---community
narratives, dream-analysis literature, and the user's personal journal---
fusing BM25 and pgvector scores with Reciprocal Rank Fusion (RRF;
\href{https://doi.org/10.1145/1571941.1572114}{Cormack et al., 2009}) and
expanding the seed set via a recursive-CTE walk over typed edges in a
Supabase \texttt{doc\_relations} table. Results are routed to roughly a
dozen widgets (interpretation synthesis, symbol graph, emotion radar,
community mirror, heatmap, stat card) that consume raw chunks, SQL
aggregates, or typed edges by modality, each carrying chunk-level
provenance. Retrieval therefore determines both \emph{what} is shown and
\emph{how} it is rendered.

\section{Background and Motivation}

\paragraph{Problem.}
Dream journaling produces rich longitudinal text, but journalers,
researchers, and clinicians lack scalable tools to connect private entries
to the accepted interpretive frameworks---Jungian symbolism, Freudian
latent-content analysis, the Hall/Van de Castle (HVdC) coding system, and
Domhoff's continuity-hypothesis account---or to the much larger corpus of
community dreams in public datasets. Current practice is manual symbol
lookup and cross-referencing that does not scale.

\paragraph{Why it is hard.}
Single-pass RAG returns a flat top-$k$ list and cannot support multi-hop
symbolic reasoning, cross-corpus triangulation, or temporal co-occurrence
discovery inside a user's archive, and it lacks the provenance required
when mixing clinical literature with personal data. The corpora also have
heterogeneous register (informal journal, HVdC codes, academic prose) and
pervasive symbolic indirection---``water'' should reach Jungian
``unconscious'' through typed edges, not lexical overlap---across three
namespaces, with every interpretation tagged back to the chunks that
support it.

\paragraph{Related literature.}
Interpretive backbone: the Hall/Van de Castle manual
(\href{https://dreams.ucsc.edu/Coding/}{dreams.ucsc.edu/Coding}), Jung,
Freud, and Domhoff (MIT Press, 2018); NLP-scale precedent:
\emph{Our Dreams, Our Selves} (2020) and the Reddit \texttt{r/Dreams}
topic-model study
(\href{https://doi.org/10.1140/epjds/s13688-025-00530-4}{Sanz et al., 2025}).
Technical foundations: RRF (Cormack et al., 2009), HNSW
(\href{https://arxiv.org/abs/1603.09320}{Malkov \& Yashunin, 2018}), and
Qwen3-Embedding
(\href{https://arxiv.org/abs/2506.05176}{Zhang et al., 2025}), whose
instruction-prefixed queries we adopt.

\paragraph{Data.}
We restrict the community corpus to sources that ship with structured
annotations so HVdC codes come from the source rather than an LLM.
Table~\ref{tab:data} lists the two ingested community sources
($\sim$48K annotated narratives), the textbook layer, and per-user
Supabase entries. Unannotated variants (DreamBank/DReAMy-lib, SDDb,
Reddit \texttt{r/Dreams}) are deferred pending a runtime auto-tagger.

\begin{table}[h]
\centering
\small
\begin{tabular}{@{}l l p{7.2cm}@{}}
\hline
Source & Scale & Annotations / notes \\
\hline
DreamBank Annotated & $\sim$28K reports & HVdC emotion + character codes (HF \texttt{gustavecortal/DreamBank-annotated}) \\
Dryad HVdC-annotated & $\sim$20K reports & Full HVdC codes + numeric indices (\href{https://datadryad.org/dataset/doi:10.5061/dryad.qbzkh18fr}{Dryad 10.5061/dryad.qbzkh18fr}) \\
Freud, Jung (Gutenberg) & textbook & Chunked at 1500/200 \\
HVdC coding manual & reference & Scraped from \href{https://dreams.ucsc.edu/Coding/}{dreams.ucsc.edu/Coding} \\
\hline
\end{tabular}
\caption{Ingested sources. Per-user Supabase entries are populated at runtime
and auto-tagged only when source annotations are absent.}
\label{tab:data}
\end{table}

\section{Proposed Approach / Implementation Details}

\begin{itemize}
  \item \textbf{Pre-processing.} Corpora are ingested through the
  \texttt{RAGStore}/\texttt{DataAdapter} components of the Supabase RAG
  scaffold. Academic texts chunk at 1500/200 tokens (1200/150 for the
  scraped HVdC manual); dream narratives are stored at report level. HVdC
  emotion, character, aggression, friendliness, and sexuality codes---plus
  numeric HVdC indices on Dryad---are parsed from source CSV/TSV columns
  into \texttt{metadata}, so bulk ingest needs no LLM tagging; a planned
  4B-class Qwen tagger will cover user journal entries without source
  annotations. Edges are extracted at ingest: in-archive symbol
  co-occurrences strengthen on repetition, and cross-corpus
  \texttt{similar\_to} edges link personal, community, and academic chunks.

  \item \textbf{Retrieval algorithm.} Retrieval runs entirely in Supabase
  Postgres through the scaffold's \emph{Context Mesh}. RRF first fuses BM25
  and pgvector cosine scores over \texttt{vector(1024)} HNSW embeddings
  into a seed set; a recursive-CTE walk over \texttt{doc\_relations} then
  expands it along seven typed edges (\texttt{symbolizes},
  \texttt{similar\_to}, \texttt{co\_occurs}, \texttt{follows},
  \texttt{interprets}, \texttt{contradicts}, \texttt{coded\_as}). The pipeline is a self-correcting LangGraph:
  \texttt{classify\_query} $\rightarrow$ \texttt{plan\_retrieval}
  $\rightarrow$ \texttt{retrieve} $\rightarrow$ \texttt{grade}
  $\rightarrow$ \texttt{rewrite\_retry} $\rightarrow$
  \texttt{synthesize}.

  \item \textbf{Widget composition.} The \texttt{synthesize} node emits
  layout JSON mapping retrieval results onto $\sim$14 widget types, each
  consuming a different slice: raw chunks (\emph{Interpretation Synthesis},
  \emph{Textbook Card}), typed edges (\emph{Symbol Network} on
  \texttt{co\_occurs}, \emph{Dream Timeline} on \texttt{follows},
  \emph{Community Mirror} on cross-corpus \texttt{similar\_to}), or SQL
  aggregates over \texttt{user\_dreams} (\emph{Heatmap}, \emph{Emotion
  Radar}, \emph{Stat Card}). Composition rules keyed on query type
  (symbol, temporal, similarity, new entry) pick the subset and grid size,
  so the same query vocabulary yields structurally different dashboards.

  \item \textbf{Robustness.} The \texttt{grade} node runs an LLM relevance
  check that triggers query rewrite on failure, using a different model
  family than the synthesizer to reduce self-preference bias. The three
  namespaces (\texttt{community\_dreams}, \texttt{dream\_knowledge},
  \texttt{user\_\{uid\}\_dreams}) are queried in isolation with
  instruction-prefixed queries. Class imbalance---$\sim$48K community
  reports vs.\ textbook layer vs.\ per-user archives (tens to
  hundreds)---is handled by stratified per-namespace budgets and HVdC
  frequency capping; graph over-fitting is controlled by normalizing
  \texttt{co\_occurs} weights by user entry count and validating on a
  held-out user split.

  \item \textbf{Inference stack and scale.} All inference is local via
  \texttt{llama.cpp} / \texttt{llama-server} on a GCP
  \texttt{g2-standard-8} instance with an NVIDIA L4 (24\,GB VRAM); no cloud
  APIs in the serving path. Chat runs
  \texttt{unsloth/Qwen3.6-35B-A3B-GGUF:UD-IQ4\_XS} (17.7\,GB)---a MoE
  activating $\sim$3B parameters per token, so throughput tracks a dense
  3B rather than a 35B. IQ4\_XS matches Q4\_K\_M quality within noise
  while saving $\sim$4.7\,GB over UD-Q4\_K\_XL; the freed headroom lets
  the KV cache stay at \emph{fp16} (not q8\_0) for 32K context under
  \texttt{parallel=1} ($\sim$3\,GB KV plus buffers/CUDA context, with
  slack for a second slot). fp16 KV matters because q8\_0 KV degrades
  long-context needle retrieval and tool-call argument fidelity---both
  load-bearing for the graph-hop RAG and agent loop. A dense Qwen3.5-27B
  @ Q4 ($\sim$16\,GB) is retained as a fallback. Embeddings use
  Qwen3-Embedding-0.6B (Q8\_0, 1024-dim MRL, last-token pooling,
  L2-normalized); MRL can drop query-time dimension to 768/512 if latency
  binds. HNSW on \texttt{vector(1024)} plus GIN indexes on symbol/emotion
  tag arrays support hybrid filtering; namespace-scoped queries bound
  the recursive CTE.

  \item \textbf{Validation.} Retrieval: Precision@5 and Recall@5 on 500
  curated DreamBank similar-dream pairs (two labelers, adjudicated) across
  four baselines---BM25, pgvector, RRF, RRF$+$graph---isolating the graph
  layer. Interpretation: blinded 5-point Likert by three raters on a
  held-out Dryad HVdC partition (excluded from graph training) with
  Cohen's $\kappa$. Provenance: manual audit on 50 cards that linked
  chunks support rendered content and graph-hop metadata is correct. User
  study: two-week within-subjects journaling ($N{=}15$) comparing a static
  dashboard against the composed layout on perceived insight,
  pre-registered for $d{=}0.5$. We also report the share of synthesis
  content from depth-$>$0 hops. Operating targets: $<$2\,s first token,
  $<$50\,ms per embedding.

  \item \textbf{Limitations.} Provenance ensures retrieval traceability
  but not generator faithfulness, so a post-synthesis faithfulness check
  runs against linked chunks. Cold-start users get seed-only layouts
  until their archive reaches a minimum size. Low-bit MoE quantization
  may degrade long-chain reasoning, probed by comparing UD-IQ4\_XS to a
  higher-bit reference. All corpora are English and Western-theory-
  dominated, limiting cross-cultural generalization. DreamRAG is a
  research prototype, not a clinical tool, and renders disclaimers on
  every interpretation card.
\end{itemize}

\end{document}
